\chapter{Introduction}
\label{ch:intro}

\section{Motivation: Physical AI meets the memory wall}

A robot that moves through the world must answer two questions continuously
and in real time: \emph{where am I?} and \emph{what does my surrounding look
like?} The computational discipline that answers both at once from camera
images is visual Simultaneous Localization And Mapping (SLAM), and in
production robotics it runs on a GPU. The workload studied in this thesis is
cuVSLAM~\cite{cuvslam2025}, NVIDIA's CUDA-accelerated visual SLAM library,
shipped inside the Isaac robotics stack and deployed on real robots. It is,
by construction, a representative workload of what is now called
\emph{Physical AI}: machines that sense and act in the physical world under
hard latency and energy budgets.

For such machines, the binding constraint on the processor is no longer
arithmetic. Four decades of technology scaling have made floating-point
operations nearly free relative to the cost of \emph{moving the operands}:
fetching a word from off-chip DRAM costs orders of magnitude more energy than
computing with it, and the gap widens every generation. This is the
\emph{memory wall}~\cite{wulf1995memorywall}, the modern face of the von
Neumann bottleneck. For a battery-powered robot the consequence is direct:
joules spent moving pixels and map entries are joules not spent on flight
time or actuation.

The architecture community's response is \emph{memory-centric} hardware:
processing-in-memory (PiM), which places compute engines at or inside the
DRAM devices~\cite{mutlu2019pim,ghose2019pim}; in-storage or near-storage
processing, which executes queries where the data is
stored~\cite{do2013smartssd,gu2016biscuit}; and near-sensor processing, which
consumes pixel streams before they ever reach main memory. Commercial silicon
exists for each of these ideas~\cite{gomezluna2021upmem,kwon2021hbmpim,
lee2022aim}. What does \emph{not} exist is a principled way to decide, for a
given production workload, \emph{which} of its computations belongs on
\emph{which} substrate. That decision requires measurement: a
\emph{characterization} of exactly how the workload moves data.

\section{Problem statement}

The problem this thesis addresses can be stated in one sentence:

\begin{quote}
\emph{For each GPU kernel of a production visual-SLAM system, determine---by
measurement, at the granularity of named data structures, with quantified
confidence---whether its bottleneck is data movement, which class of data
movement, and consequently which hardware substrate (GPU, CPU, DRAM-PiM,
scatter-capable PiM, near-sensor, or in/near-storage) is the appropriate home
for it.}
\end{quote}

Three properties make this problem hard, and each drives a contribution of
this thesis:

\begin{enumerate}
\item \textbf{The reference methodology is CPU-shaped.} The established
data-movement characterization methodology, DAMOV~\cite{oliveira2021damov},
was built for CPUs. GPUs violate its assumptions in specific ways: they have
a two-level (not three-level) cache hierarchy; they hide memory latency with
massive thread parallelism rather than stalling; and they possess a failure
mode with no CPU analogue---\emph{memory divergence}, where the 32 threads of
a warp scatter their accesses and waste up to 8$\times$ the fetched
bandwidth. A GPU characterization must adapt the methodology, not merely
reuse it (Chapter~\ref{ch:method}).
\item \textbf{Kernel-level claims are not architect-level claims.} A profiler
can report that a kernel is memory-bound; an architect needs to know
\emph{which data structure} the traffic belongs to---an image, a solver
matrix, or a growing keyframe database---because the substrate verdict
differs for each. Attributing every byte of DRAM traffic to a named
allocation in a closed-source production binary requires purpose-built
instrumentation (Chapter~\ref{ch:attribution}).
\item \textbf{Measurements must be proven trustworthy.} A hostile reviewer
will ask, in order: Was the workload even computing correctly? Did the
instrumentation change what was measured? Are the numbers stable run to run?
Would the classification survive different thresholds, different clustering
algorithms, different hardware? Each of these questions is answered with a
dedicated experiment in this thesis (Chapters~\ref{ch:validity}
and~\ref{ch:characterization}).
\end{enumerate}

\section{Research questions}

\begin{description}
\item[RQ1] Can DAMOV's data-movement methodology be re-derived for GPUs such
that each of its components---screening, machine-independent locality
analysis, bottleneck classification, and robustness validation---has a
measured GPU analogue? (Chapter~\ref{ch:method};
validated in Chapters~\ref{ch:characterization}--\ref{ch:locality}.)
\item[RQ2] What are the memory-bottleneck classes of a production visual-SLAM
workload, and are they properties of the kernels themselves or artifacts of
particular inputs, thresholds, or machines?
(Chapter~\ref{ch:characterization}.)
\item[RQ3] Which named data structures receive each kernel's DRAM traffic,
and what fraction of that traffic is program data versus compiler-generated
scratch? (Chapter~\ref{ch:attribution}.)
\item[RQ4] What is the resulting substrate placement for the workload's
computation, what fraction of GPU time is offload-eligible, and what
measured baselines (energy, host I/O) must a follow-on architecture study
improve upon? (Chapter~\ref{ch:synthesis}.)
\item[RQ5] Can the entire measurement apparatus be shown to be
accuracy-neutral---i.e., that the characterized system is a correctly
functioning SLAM whose outputs are unchanged by observation?
(Chapter~\ref{ch:validity}.)
\end{description}

\section{Contributions}

\begin{enumerate}
\item \textbf{GPU-DAMOV, a methodology} (Chapter~\ref{ch:method}). A
component-by-component re-derivation of DAMOV for GPUs: LFMR redefined for a
two-level hierarchy; latency-boundedness conditioned on occupancy; a
coalescing axis added; DAMOV's simulated intervention experiment replaced by
two \emph{measured} interventions (a clock-domain sensitivity sweep and a
direct reuse-distance cache-capacity curve) with the simulated third axis
prepared as generated configuration. The method ends one step later than
DAMOV: not at ``memory-bound'' but at a per-kernel \emph{substrate verdict}
plus, when the verdict is ``stay on GPU,'' the named architectural fault to
fix.
\item \textbf{The first per-data-structure memory characterization of a
production SLAM system} (Chapters~\ref{ch:characterization}--\ref{ch:synthesis}).
49 kernels $\times$ 27 sequences $\times$ 4 datasets $\times$ 192
configuration variants, with an eight-class taxonomy (G0--G7) in which class
membership is 91\% consistent across sequences, independently recovered by
two clustering algorithms, and 80\% reproducible across two GPU
microarchitectures.
\item \textbf{TaggedAllocator attribution} (Chapter~\ref{ch:attribution}). A
three-layer instrumentation---source-level allocation journaling with call
stacks, driver-level allocation lifetime capture, and a streaming join
against binary-instrumentation address traces---that resolves 274/274
allocations and gives 48/49 kernels a unanimous data-structure tag. It
surfaces a finding invisible to counters: 92.6\% of the loop-closure scan's
DRAM volume is register-spill scratch, not data.
\item \textbf{A validated trust chain} (Chapter~\ref{ch:validity}). A
141-configuration accuracy matrix reproducing the vendor's published
trajectory accuracy; a 192-configuration campaign showing profiling is
accuracy-neutral (bit-identical trajectories on deterministic modes); a
measured noise floor (run-to-run CoV 0.14\% at locked clocks); and a
$\pm$25\% threshold sensitivity analysis.
\item \textbf{A reusable, workload-agnostic toolchain} (released under MIT
license): a single-file experiment runner, a config-mutation regime, a
three-profiler harness with adapters for arbitrary GPU workloads, a
stdlib-only analysis library that regenerates every table from committed
CSVs without a GPU, and an interactive evidence dashboard.
\item \textbf{Two documented self-corrections} in which independent methods
disagreed and were reconciled on record---one reversing an earlier
project-internal conclusion (Chapter~\ref{ch:locality})---presented as a
feature of the methodology: the pipeline catches its own errors.
\end{enumerate}

\section{Scope and non-goals}

This thesis is a \emph{characterization}, deliberately stopping at measured
candidacy plus an analytical placement bound. It does \emph{not} simulate a
PiM or in-storage design and therefore reports no end-to-end speedup or
energy delta for the proposed substrates; that is the follow-on architecture
study, for which this thesis generates the simulator configurations and
measures the baselines (Section~\ref{sec:roadmap}). It characterizes one
workload family---cuVSLAM---arguing representativeness from production
deployment rather than breadth across SLAM implementations; the toolchain's
adapter mechanism is the demonstrated path to breadth. Hardware evaluation
uses one desktop-class GPU as the primary instrument, with a second
microarchitecture used for cross-device validation of the classification.

\section{Thesis organization}

Chapter~\ref{ch:background} builds the technical background from first
principles: the GPU execution model, the memory hierarchy and its five
address spaces, SLAM and cuVSLAM, the candidate substrates, and the
measurement instruments. Chapter~\ref{ch:related} situates the work among
data-movement characterizations, PiM systems, and SLAM acceleration studies.
Chapter~\ref{ch:method} presents GPU-DAMOV: metrics, decision tree,
taxonomy, placement model, and the experimental setup.
Chapter~\ref{ch:validity} establishes trust: accuracy validation, profiler
neutrality, and measurement rigor. Chapter~\ref{ch:characterization} reports
the kernel-level characterization and its four-way validation.
Chapter~\ref{ch:locality} reports the machine-independent locality analysis
and the first self-correction. Chapter~\ref{ch:attribution} reports the
data-structure attribution and the register-spill finding.
Chapter~\ref{ch:synthesis} synthesizes the substrate verdicts, the placement
model, and the measured energy and host-I/O baselines.
Chapter~\ref{ch:threats} enumerates threats to validity.
Chapter~\ref{ch:conclusion} concludes with the phased roadmap. Appendices
provide a glossary, the artifact/reproducibility description, and extended
result tables.
